%%%%%%%%%%%%%%%%%%%%%%%%%%%%%%%%%%%%%%%%%%%%%%%%%%%%%%%%%%%%%%%%%%
%% Toto je inkludovaný LaTeXový soubor `x_levit.tex'.           %%
%% This is the `x_levit.tex' included LaTeX file.               %%
%% ------------------------------------------------------------ %%
%% Verze / Version 2022-05-23                                   %%
%% Vyvíjí & udržuje / Developed & maintained by                 %%
%%                           <stanislav.hledik@physics.slu.cz>  %%
%%%%%%%%%%%%%%%%%%%%%%%%%%%%%%%%%%%%%%%%%%%%%%%%%%%%%%%%%%%%%%%%%%

\chapter{Rotací stabilizovaná magnetická levitace}\label{SpinLev}%%%%%%%

Pokud jste se někdy snažili přimět jeden permanentní magnet k~levitaci nad
jiným permanentním magnetem (kdo toto aspoň jednou nezkoušel?), jistě jste
zjistili, že lze snadno zajistit mezi magnety repulsivní sílu dostatečně
silnou pro vybalancování tíhy. Stejně jste ale asi k~vaší nelibosti poznali,
že levitátor~-- magnet, který měl levitovat, se vždy překlopil, čímž se
magnetická repulse změnila na atrakci, levitátor spadl a \uv{přilepil} se na
druhý magnet. Ať jste se snažili sebevíc, stabilní levitace nebylo možno
dosáhnout. Nemožnost stabilní levitace permanentních magnetů teoreticky
odvodil a formalizoval Reverend \designation{Samuel Earnshaw} v~roce 1842
(tedy ještě před publikací Maxwellových rovnic klasické elektrodynamiky roku
1861) kdy publikoval svůj slavný
teorém~\cite{Ear:1842:TCAMPS:NatMoFoConsLumEth,%
  Bas:2006:MECJI:EarnshawMagLev}, jemuž se budeme věnovat podrobně
v~následující sekci~\ref{EarnshawTheorem} na
straně~\pageref{EarnshawTheorem}.\footnote{Earnshawova motivace byla ovšem
  jiná~-- týkala se nestability rovnováhy interagujících elektrických nábojů a
  její aplikace na světlonosný éter. Platnost Earnshawova teorém je ovšem
  možné snadno rozšířit na interakci libovolných elektrických či magnetických
  multipólů.} Od té doby považovali fyzikové levitaci permanentních magnetů za
nemožnou, analogicky jako matematici považují (oprávněně) za nemožné roztřetit
pomocí pravítka a kružítka libovolný úhel (což bylo dokázáno rovněž
v~19.~století).

Ukazuje se, že Earnshawův teorém je sice správný, ale v~jeho
předpokladech jsou \uv{zadní vrátka} umožňující přece jen teoreticky
zdůvodnit a realizovat dokonce více druhů magnetické levitace. Některé z~nich
jsme ve stručnosti představili v~předcházející kapitole. Nyní se budeme
věnovat obzvláště fascinující magnetické levitaci, jež nevyžaduje napájení,
kryogenní aparaturu, a ani žádné speciální prostředky. Je tak jednoduchá, že
je komerčně k~mání jako \uv{hračka} či \uv{hlavolam} za cca USD~40 pod názvem
\Levit{} (viz~\cite{Fas:LIVE:Levit:Levitron,%
  Gra:LIVE:Levit:IgnMagBliss}, lze se setkat i~s~pojmenováním
U-CAS~\cite{Gov-Sht:1999:phys9902002:HowHiUCAS}).

Symetrický vlček (levitátor) obsahující permanentní magnet rotující a
vznášející se nad bází obsahující toroidální permanentní magnet si nechal
v~roce 1983 patentovat \designation{Roy Harrigan} z~Vermontu,
USA~\cite{USpatent4.382.245}. Tento vynálezce měl vůči svým předchůdcům, kteří
se před ním snažili o~levitaci permanentních magnetů (a~neuspěli), značnou
výhodu: neznal Earnshawův teorém. Protože nevěděl, že permanentní magnety
nelze levitovat, náhodou a usilovným snažením objevil způsob, jak toho
dosáhnout \emph{lze}. To není protimluv~-- Harrigan nevědomky využil toho, že
rotace levitátoru na sebe váže dva jeho stupně volnosti, které se tak nemohou
měnit nezávisle, nýbrž změna jednoho velmi rafinovaným způsobem ovlivňuje
druhý. Tím vytváří ostrůvek stability způsobem, který nenarušuje Earnshawův
teorém, jenž předpokládá nezávislost všech stupňů volnosti levitátoru. Takovou
\uv{změnu předpokladů} nikdo z~renomovaných fyziků, kteří Harrigana
přesvědčovali o~marnosti jeho úsilí, po více než století od formulace
Earnshawova teorému nenavrhl~\cite{Gra:LIVE:Levit:IgnMagBliss}.\footnote{Je
  dobré si připomenout tuto historii kdykoli Vám někdo řekne, že něco
  nejde. Nejsou-li k~tomu racionální důvody, je lepší předpokládat, že to
  udělat jde.} Nutno ovšem podotknout, že vysvětlení fyziky levitace
permanentních magnetů je na mnoha webových stránkách jen částečně správně,
včetně stránky~\cite{Fas:LIVE:Levit:Levitron}.

Tato práce čerpá převážně ze dvou základních článků a jedné diplomové
práce~\cite{Sim-Hef-Rid:1997:AMEJP:SpStMagLev,%
  Ber:1996:PRORSA:LevitAdiaTrapSp,%
  Hab:1999:MgrDP:ProbUlFyzPrakt},
částečně z~dalších seriózních
zdrojů~\cite{Caz-Nem:2010:REVRS:NovConPerMagLev,%
  Dul-Eas:1999:ZANGM1:StabOfLevit,%
  Fla-etal:1999:PHYSD:SpiMoHovMaTop,%
  Fre:2010:PHYSED:WondersLevit,%
  Gan-Jon-Was:1998:JPHYSD:DynOfLevi,%
  Gen-Del-Ron:1999:MECJI:GyStaPasMagLev,%
  Gov-Sht-Tho:1999:PHYSD:DyStHoMagTop,%
  Gov-Sht:1999:phys9902002:HowHiUCAS,%
  Jon-Was-Gans:1997:JAPPPH:SimpTheoLevi}
i~z~dalších, zejména webových
zdrojů~\cite{Ber:LIVE:Levit:LevitronPhys,%
  Bor:LIVE:Levit:OptRotLevi}. K~práci~\cite{Gan-Jon-Was:1998:JPHYSD:DynOfLevi}
je na webu zdrojový soubor \pgm{dyn\_of\_lev.mws} pro Maple numerické simulace
Hanse-Georga Koepkena.

\begin{figure}[p]
\begin{minipage}[b]{0.627\linewidth}
\begin{center}

\includegraphics[width=\linewidth]{Levi1}
\end{center}
\end{minipage}\hfill%
\begin{minipage}[b]{0.333\linewidth}
\begin{center}
\includegraphics[width=\linewidth]{Levi2}
\end{center}
\end{minipage}
\caption[Schéma Levitronu]{\label{Levi}\emph{Vlevo:} Průřezové schéma
  Levitronu: (1)~horní kónická část levitátoru ve formě symetrického
  setrvačníku (vlčku) určená k~uchopení a roztáčení; (2)~kruhový keramický
  permanentní magnet; (3)~pomocná plastová zdvihací podložka; (4)~plastové
  pouzdro bázového magnetu; (5)~toroidální keramický permanentní bázový
  magnet. \emph{Vpravo:} Prostorové schéma Levitronu s~vyznačenými
  inerciálními kartézskými souřadnicemi se středem v~centru bázového magnetu,
  vektorem magnetického dipólového momentu $\vec{\mu}$, lokálním vektorem
  magnetické indukce $\vec{B}$, normálovým vektorem $\vec{S}$ úhlové rychlosti
  ($|\vec{S}|=\omega$) a těžištěm levitátoru~G. Převzato
  z~\cite{Gen-Del-Ron:1999:MECJI:GyStaPasMagLev}.}
\end{figure}
\begin{figure}[p]
\begin{minipage}[b]{0.61\linewidth}
\begin{center}
\includegraphics[width=\linewidth]{Levi3}
\end{center}
\caption[Realizace Levitronu]{\label{Levi3}Improvizovaná realizace
  levitátoru Levitron.}
\end{minipage}\hfill%
\begin{minipage}[b]{0.35\linewidth}
\begin{center}

\includegraphics[width=\linewidth]{Earnshaw}
\end{center}
\caption[Samuel Earnshaw]{\label{Earnshaw}Samuel
  Earnshaw. }
\end{minipage}
\end{figure}

\Levit{} a jeho koncepční schéma lze vidět na
%frontispice (obrázek převzat z~\cite{Fas:LIVE:Levit:Levitron}) a na
obrázcích~\ref{Levi} a~\ref{Levi3}. Pokud levitátor umístíme do vhodné pozice
nad toroidní magnet a roztočíme jej s~úhlovou rychlostí nacházející se
v~určitých mezích, dojde k~jeho stabilní levitaci, která bude trvat do doby,
než se levitátor vlivem tření zpomalí pod dolní mez úhlové
rychlosti. Přivedení levitátoru do režimu levitace vyžaduje určitý
cvik. Postup je následující: bázový magnet umístíme na stolek, přičemž dbáme
na přesně horizontální nastavení. Položíme na něj plastovou destičku
(dodávanou se zařízením \Levit{}), do jejíhož středu umístíme levitátor a
prsty jej roztočíme. Následně pomalu zdvihneme destičku do výše asi
2.5--3\,cm, až ucítíme, že levitátor v~blízkosti levitační výšky jakoby
z~destičky \uv{vyskočil}.  Nyní můžeme již destičku odstranit a sledovat
levitaci. V~případě, že levitátor má snahu \uv{utéci} do strany, musíme
příslušnou stranu pomocí stavitelných nožiček stolku zdvihnout. Během zdvihání
plexisklové destičky se levitátor stabilizuje v~horizontálním směru, což
v~usnadňuje jeho vyzdvižení~-- důvod je vysvětlen v~poznámce pod čarou za
rovnicí~(\ref{dUXYlt0}). Je proto dobré co nejdříve po roztočení levitátoru
nadzvednout plexisklovou destičku do výše asi 0.5~cm, kde dochází ke
stabilizaci a snižuje se tak precese. Při správném vyvážení stolku a vhodné
rotaci je levitátor schopen levitovat po dobu asi 3~minut, než se
aerodynamickým třením sníží otáčky natolik, že gyroskopická stabilizace
přestane fungovat.

Principy rotačně stabilizované levitace permanentních magnetů se ovšem úzce
prolínají s~principy diamagnetické levitace, jak se čtenář může přesvědčit
z~výše uvedené literatury.

Uveďme zde několik odkazů na videoklipy presentující \Levit{} nebo jeho
napodobeniny v~akci:
\begin{itemize}
\item\url{https://youtu.be/YpHraxs6Le4}
\item\url{https://youtu.be/tvDRCwIS8KY}
\item\url{https://youtu.be/GMVtlNbMwHw}
\item\url{https://youtu.be/JOZeCTF_Ilk}
\item\url{https://youtu.be/VylZJrep1MM}
\item\url{https://youtu.be/9LTENE7TYew}
\item\url{https://youtu.be/iRZfk_FvS98}
\item\url{https://youtu.be/1yVko_ANnMc}
\end{itemize}

\section{Earnshawův teorém}\label{EarnshawTheorem}%/////////////////////

Tento teorém, který roku 1842 odvodil \designation{Samuel Earnshaw} zmíněný
v~úvo\-du kapitoly (viz obrázek~\ref{Earnshaw}), má v~problematice magnetické
levitace zásadní význam. Uveďme nejprve několik jeho slovních formulací.

\begin{quote}\itshape \label{ETorig}
  Jakákoli rovnovážná statická konfigurace elektrických nábojů, fixních
  dipólů, kvadrupólů nebo vyšších multipólů umístěných volně v~prostoru je
  nestabilní.~[Formulace blízká originální]
\end{quote}
Totéž ovšem musí platit i~pro permanentní magnety (fixní magnetické dipóly
nebo multipóly vyšších řádů), neboť vztahy pro interakční energii elektrických
a magnetických multipólů jsou stejné (s~výjimkou neexistujících magnetických
monopólů).\footnote{Známá je úloha ze středoškolské fyziky o~nalezení
  elektrického náboje $Q$, který je třeba umístit do středu čtverce, v~jehož
  vrcholech jsou stejné elektrické náboje $q$, tak, aby celková síla na každý
  z~těchto pěti nábojů vymizela. Nalézt takovou rovnováhu je triviální,
  stabilita je však nedosažitelná~-- sebemenší perturbace polohy kteréhokoli
  z~nábojů vede ke kolapsu systému.} Dodatečné vnější gravitační pole výsledky
nemění, neboť jeho relevantní vlastnosti jsou stejné jako
u~elektrostatického či magnetostatického pole (viz níže).

\begin{quote}\itshape \label{ETshg}
  Neexistuje čistě elektrostatický nebo magnetostatický levitátor nebo
  částicová past.~\cite{Sim-Hef-Gei:2001:AMEJP:DiaStMagLev}
\end{quote}

\begin{quote}\itshape \label{ETsgbg}
  Žádná kombinace elektrostatických, magnetostatických nebo gravitačních
  sil nemůže vytvořit 3D potenciálovou jámu nezbytnou pro stabilní levitaci ve
  volném
  prostoru.~\cite{Sim-Gei:2000:JAPPPH:DiaLevFlyFrog,%
    Ber-Gei:1997:EURJP:OfFlyFrogLev}
\end{quote}

\begin{quote}\itshape \label{ETgetal}
  Žádný stacionární objekt vyrobený z~permanentních magnetů ve fixní
  konfiguraci nelze udržet ve stabilní rovnováze jakoukoli kombinací
  statických magnetických a/nebo gravitačních
  sil.~\cite{Gei-etal:1999:NATURE:MagLevFinger}
\end{quote}

Dokážeme nyní Earnshawův teorém pro libovolný systém magnetických dipólů,
případně i~multipólů vyšších řádů, jenž se mohou v~prostoru libovolně
přemisťovat a natáčet, a to způsobem, že všechny jejich stupně volnosti jsou
navzájem nezávislé. Navíc předpokládejme, že magnetické multipólové momenty
mají také hmotnost nebo analogicky momenty vyšších řádů popisující rozložení
hmoty, takže podléhají gravitačnímu působení.\footnote{Přidání elektrických
  multipólů je analogické, ale protože je nebudeme v~dalších úvahách
  potřebovat, v~zájmu zjednodušení je vypustíme.} Uvažujme takový systém $N$
vzájemně interagujících magnetických a hmotných multipólů očíslovaných
indexem $a=1,\ldots,N$; mezi nimi předpokládáme vakuum nebo prostředí
homogenní hustotou hmoty i~magnetickými vlastnostmi. Vybereme jeden
multipól~-- nechť je to bez újmy na obecnosti multipól~1, jehož polohu
symbolicky označíme symbolem $P$. Magnetickou indukci, gravitační intenzitu a
potenciál vytvořené v~bodě $P$ všemi ostatními multipóly $a=2,\ldots,N$
označme po řadě $\vec{B}(P)$, $\vec{g}(P)$ a $V(P)$. Bod $P$ se nachází mimo
ostatní multipóly, takže pro gravitační potenciál platí Laplaceova
rovnice~\cite{Gol-Poo-Saf:2002:ClassMech:3}
\begin{equation}
  \Delta V(P) = 0\,.                                             \label{harmonV}
\end{equation}
Jelikož tato rovnice platí i~v~určitém okolí bodu $P$, v~němž se nechacházejí
ostatní členové systému, budou v~tomto okolí nulové i~parciální derivace
gravitačního potenciálu podle kartézských souřadnic libovolného řádu, a po
záměně pořadí parciálních derivací a s~přihlédnutím ke vztahu mezi gravitačním
potenciálem a intenzitou $\vec{g} = -\nabla V$ dostaneme
\begin{equation}
  \Delta g_{i} = \Delta\pder{g_{i}}{x_{j}}
    = \Delta\frac{\partial^{2}g_{i}}{\partial x_{j}\partial x_{k}} = \cdots
    = 0\,,                                                       \label{harmonG}
\end{equation}
kde $i=1,2,3$ a $(x_{1},x_{2},x_{3})\equiv (x,y,z)$ jsou kartézské inerciální
souřadnice a $(g_{1},g_{2},g_{3}) = (g_{x},g_{y},g_{z})$ odpovídající
komponenty gravitační intenzity. Rovnice analogická~(\ref{harmonG}) platí
i~pro komponenty magnetické indukce,
\begin{equation}
  \Delta B_{i} = \Delta\pder{B_{i}}{x_{j}}
    = \Delta\frac{\partial^{2}B_{i}}{\partial x_{j}\partial x_{k}} = \cdots
    = 0\,.                                                       \label{harmonB}
\end{equation}
Důkaz je jednoduchý a ihned plyne z~již dříve zmíněné nezřídlovosti
($\nabla\cdot\vec{B}=0$) a nevírovosti magnetické indukce
($\nabla\times\vec{B}=\vec{0}$)~\cite{Hab:1999:MgrDP:ProbUlFyzPrakt} pomocí
známé identity~\cite{Pur-Mor:2013:ElectMagn:3}
\begin{equation}
  \Delta\vec{B}
  =\nabla(\nabla\cdot\vec{B})-\nabla\times(\nabla\times\vec{B})\,.\label{Errgdl}
\end{equation}
Pole, jejichž laplacián je nulový, nazýváme harmonické. Podle
rovnic~(\ref{harmonV}), (\ref{harmonG}) a (\ref{harmonB}) jsou gravitační
potenciál, komponenty gravitační intenzity a magnetické indukce včetně všech
jejich parciálních derivací harmonická pole.

Potenciální energie vybraného multipólu je~dána
vztahem~\cite{Pur-Mor:2013:ElectMagn:3}
\begin{align}
  U(P,Q) &= mV(P) + \left(\text{hmotový dipólový, kvadrupólový
   a vyšší členy}\right)\notag\\
         &- \vec{\mu}(Q)\cdot\vec{B}(P)
   + \left(\text{magnetický kvadrupólový a vyšší členy}\right)\,,  \label{PotEn}
\end{align}
v~němž $m$ označuje hmotnost, $\vec{\mu}$ magnetický dipólový moment a $Q$
symbolicky vnitřní stupně volnosti (u~dipólu jeho natočení např. pomocí
Eulerových úhlů). Členy popisující vyšší multipóly jsou vyznačeny pouze
symbolicky, je však jasné, že budou multiplikativním způsobem obsahovat
parciální derivace gravitačního potenciálu a komponent $B_{i}$ podle
prostorových souřadnic~\cite{Lan-Lif:1980:2-ClasTheField:4}. Z~těchto úvah
a rovnic~(\ref{harmonV}), (\ref{harmonG}) a (\ref{harmonB}) dostáváme, že
potenciální energie vybraného multipólu \uv{dědí} vlastnost harmonicity,
\begin{equation}
  \Delta U(P,Q)\equiv
  \pdder{U}{x} + \pdder{U}{y} + \pdder{U}{z} = 0\,.              \label{harmonU}
\end{equation}

V~dalším textu se omezíme na magnetické dipóly popsané vektorem magnetického
dipólového momentu $\vec{\mu}$ a hmotové monopóly $m$. Započtení multipólů
vyšších řádů je pouze technickou záležitostí a na principu nic nemění, navíc
vzhledem k~obvyklé homogenitě gravitačního pole je zanedbání hmotových
multipólů vyšších řádů exaktní.  Nutná podmínka rovnováhy vybraného multipólu
je, aby síla na něj působící v~bodě $P$ vymizela\footnote{Díky nevírovosti
  pole $\vec{B}$ platí $\lpder{B_{i}}{x_{j}}=\lpder{B_{j}}{x_{i}}$, a tedy
  $-\nabla(-\vec{\mu}\cdot\vec{B}) = (\vec{\mu}\cdot\nabla)\vec{B}$, jelikož
  $\mu(Q)$ závisí pouze na vnitřních stupních volnosti.}
\begin{equation}
  \vec{F}(P) \equiv -\nabla U(P,Q) =
  \underbrace{m\vec{g}}_{\text{tíha}}
  + \underbrace{(\vec{\mu}\cdot\nabla)\vec{B}}_{\text{magnet. síla}}
  = \vec{0}\,.                                                    \label{equlib}
\end{equation}
Pro dosažení stability je třeba zajistit konvexnost potenciální energie
v~libovolném směru, nutná podmínka čehož je
\begin{equation}
  \Delta U(P,Q)
  \equiv \pdder{U}{x} + \pdder{U}{y} + \pdder{U}{z} > 0\,,         \label{stabU}
\end{equation}
což je ve sporu s~podmínkou~(\ref{harmonU}). Díky harmonicitě potenciální
energie nemůže být podmínka stability~(\ref{stabU}) splněna. Pokud by
například parciální derivace $\partial^{2} U/\partial x^{2}$ a
$\partial^{2} U/\partial y^{2}$ byly kladné, bude pro zachování nulového
součtu v~(\ref{harmonU}) muset být $\partial^{2} U/\partial z^{2} < 0$, což
znamená labilitu v $z$-ovém směru. Může také nastat situace, že jedna
parciální derivace bude kladná a dvě záporné, takže nestabilní budou dva
stupně volnosti.

Tyto úvahy lze shrnout takto: potenciální energie vybraného multipólu
v~gravitačním a magnetickém poli ostatních multipólů je harmonické pole a nemá
tudíž prostorové lokální minimum. Může mít pouze sedlové body, v~nichž lze
stability dosáhnout nejvýše ve dvou ortogonálních směrech, ostatní aspoň jeden
směr musí být vždy nevyhnutelně labilní.

Situaci komplikují ještě vnitřní stupně volnosti $Q$, které jsme zatím nebrali
do úvahy, a které např. v~případě magnetického dipólu popisují jeho
prostorovou orientaci. Jak je patrno z~vlastností skalárního součinu, že
minimum potenciální energie $-\vec{\mu}\cdot\vec{B}$ (při zafixované pozici
$P$) nastává při paralelní orientaci dipólového momentu vůči magnetické
indukci, maximum při antiparalelní pozici. Na magnetický dipólový moment
$\vec{\mu}$ totiž působí kromě magnetické síly
\begin{equation}
  \vec{F} = (\vec{\mu}\cdot\nabla)\vec{B}                           \label{Fdip}
\end{equation}
také moment síly
\begin{equation}
  \vec{K} = \vec{\mu}\times\vec{B}                                  \label{Kdip}
\end{equation}
(odvození je uvedeno např. v~\cite{Pur-Mor:2013:ElectMagn:3}). Za povšimnutí
stojí, že pro vznik magnetické síly musí být pole magnetické indukce
nehomogenní, silový moment však může působit i~v~homogenním magnetickém poli.

V~dalším uvidíme, že právě antiparalelní orientace je nutná pro vznik
rovnováhy, jejíž stabilita je tak ohrožena nejen výše popsanou translační
labilitou, ale i~labilitou vůči orientaci (překlopení). Na druhé straně právě
vazba vhodných parametrů orientace dipólu na jeho pozici v~magnetickém poli
pomůže vymanit se z~podmínek Earnshawova teorému (vzájemná nezávislost
všech stupňů volnosti) a způsobit nahrazení znaménka rovnosti
v~(\ref{harmonU}) znaménkem \uv{větší než}, což v~principu umožní stabilní
levitaci.

\section{Jak dosáhnout prostorově stabilní levitace?}\label{BeyondE}%///

Nezbytnou podmínkou pro stabilitu levitace je kladný laplacián potenciální
energie~(\ref{PotEn}) v~rovnovážném bodě. První člen popisující gravitační
interakci bude vždy harmonický, jedinou možností, jak situaci změnit, je
ovlivnění druhého členu popisujícího magnetickou interakci.

Nejprve jednoduchým výpočtem ověříme, že pro kvadrát velikosti magnetické
indukce
\begin{equation}
  B^{2} \equiv \vec{B}\cdot\vec{B} = B_{x}^{2} + B_{y}^{2} + B_{z}^{2}
\end{equation}
platí~\cite{Ber-Gei:1997:EURJP:OfFlyFrogLev}
\begin{equation}
  \Delta (B^{2})
    = 2\left(\left|\nabla B_{x}\right|^{2}
           + \left|\nabla B_{y}\right|^{2}
           + \left|\nabla B_{z}\right|^{2}\right) \geq 0\,.      \label{LapB2gen}
\end{equation}
Levitace ovšem vyžaduje nehomogenním magnetické pole, v~němž musí být gradient
alespoň jedné z~komponent magnetické indukce nenulový, takže
v~(\ref{LapB2gen}) nastává ostrá nerovnost
\begin{equation}
  \Delta (B^{2}) > 0\,.                                            \label{LapB2}
\end{equation}
V~důsledku rostoucí monotonicity druhé odmocniny na oboru kladných reálných
čísel pak platí také
\begin{equation}
  \Delta B > 0\,.                                                  \label{LapB}
\end{equation}

Vztahy~(\ref{LapB2}) a (\ref{LapB}), dobře známé konstruktérům magnetických
částicových pastí, jsou esenciální ingrediencí pro možnost
levitace. Pokud totiž vezmeme místo fixního magnetického dipólu \uv{měkký}
dipól, jehož dipólový moment je úměrný vnější magnetické indukci, ale míří
opačným směrem,
\begin{equation}
  \vec{\mu} = -\kappa\vec{B}\,,                             \label{muPropMinusB}
\end{equation}
kde $\kappa$ je kladná konstanta (tedy přesně chování diamagnetik), bude
potenciální energie
\begin{equation}
  U = mV +\kappa B^{2}\,.                                        \label{diamagU}
\end{equation}
Aplikujeme-li na obě strany této rovnice Laplaceův operátor, dostaneme
podle rovnice~(\ref{LapB2}) okamžitě kýžený výsledek
\begin{equation}
  \Delta U > 0\,.                                               \label{LapUGt0}
\end{equation}
Této možnosti využívá diamagnetická levitace. Všimněme si, že odpadá nutnost
zabývat se vnitřními stupni volnosti $Q$, protože díky diamagnetickému jevu se
magnetizace materiálu a tím i~magnetický dipólový moment $\vec{\mu}$
automaticky nastavuje do správného (antiparalelního) směru, jak je patrné
z~(\ref{muPropMinusB}). Konstanta $\kappa$ je ovšem velmi malá, takže
k~realizaci diamagnetické levitace je zapotřebí velmi silného magnetického
pole.

Druhá možnost je ponechat velikost magnetického dipólového momentu konstantní,
ale nějak zařídit, aby se natáčel antiparalelně vektoru externí magnetické
indukce, formálně vyjádřeno
\begin{equation}
  \vec{\mu} = -\kappa\,\frac{\vec{B}}{B}\,,                   \label{muPropDirB}
\end{equation}
což odpovídá chování permanentního magnetu. Potenciální energie potom bude
\begin{equation}
  U = mV +\kappa B\,.                                             \label{permaU}
\end{equation}
Analogicky s~předchozím případem, po aplikaci Laplaceova operátoru na obě
strany této rovnice obdržíme podle~(\ref{LapB}) opět vztah~(\ref{LapUGt0})
připouštějící stabilitu.

Jak uvidíme v~dalším rozboru, striktní antiparalelizace magnetického dipólového
momentu $\vec{\mu}$ vůči externí magnetické indukci $\vec{B}$ není nutný,
stačí, když úhel mezi $\vec{\mu}$ a $-\vec{B}$ nepřekročí jistý mezní malý
úhel. Na rozdíl od diamagnetické levitace je zde jeden zásadní problém, a~to
jak udržet magnetický dipól v~blízkosti antiparalelní pozice, když na něj
působí silový moment $\vec{K} = \vec{\mu}\times\vec{B}$, který se snaží dipól
otočit do paralelní pozice. Jinými slovy řečeno: máme aspoň teoreticky vyřešen
problém prostorové stabilizace (stupňů volnosti $P$), ale nikoli vnitřních
stupňů volnosti $Q$.

Pro možnost stabilizace levitace permanentních magnetů tedy vyvstávají
dvě klíčové otázky:
\begin{enumerate}
\item Jak zařídit udržování antiparalelní vzájemné orientace magnetického
  dipólového momentu $\vec{\mu}$ a $\vec{B}$ (nemusí být zcela striktní ve
  smyslu poznámky výše)? Udržování antiparalelní orientace umožní vznik 3D
  potenciálové jamky okolo rovnovážné polohy.
\item Jak stabilizovat vnitřní stupně volnosti $Q$, tj. jak zařídit, aby se
  dipól vlivem momentu~(\ref{Kdip}) okamžitě nepřeklopil?
\end{enumerate}

\noindent
Jak uvidíme dále, odpovědí na obě otázky je \emph{gyroskopická stabilizace}
pomocí rotace. Je ale dobré si uvědomit, že jsme zatím dokázali pouze, že
levitace není vyloučena, a to rovnicí~(\ref{LapUGt0}), která stále připouští
konkávnost $U(P)$ v~některém směru,
\begin{equation}
\pdder{U^{2}}{x}>0\,,\quad
\pdder{U^{2}}{y}>0\,,\quad
\pdder{U^{2}}{z}>0\,.                                              \label{E4re}
\end{equation}
Diskuse a důkaz, že pro konfiguraci použité v~zařízení \Levit{} to lze, je
obsahem následující sekce.

\section{\Levit{} a teorie jeho stability}\label{LeviStab}%//////////////////////

Vyjdeme z~konfigurace typické pro \Levit{} znázorněné na
obrázku~\ref{UkazSchemaLevitronu}. Systém je tvořen prstencovým permanentním
magnetem (bází) a levitátorem~-- symetrickým vlčkem s~magnetickým dipólem,
jenž je rovnoběžný s~jeho osou symetrie. Permanentní prstencový magnet má
poloměr $a$ a jeho magnetický moment v~nepohyblivém laboratorním inerciálním
souřadném systému $(X,Y,Z)$ má komponenty $\vec{M}=(0,0,M)=M\hat{e}_{Z}$, kde
$M>0$. Osa prstence je shodná s~osou $Z$.

\begin{figure}[t]
\begin{center}

\includegraphics[width=\linewidth]{UkazSchemaLevitronu}
\end{center}
\caption[Schéma modelu Levitronu]{\label{UkazSchemaLevitronu}Schéma modelu
  Levitronu použité v~pozdější numerické simulaci.}
\end{figure}

Levitátor je charakterizován Amp\`erovým magnetickým dipólovým momentem
$\vec{\mu}$ s~komponentami $\vec{\mu}=(0,0,\mu)=\mu\hat{e}_{z}$ v~souřadném
systému $(x,y,z)$ pevně spojeném s~levitátorem a mající počátek v~jeho hmotném
středu\footnote{Kvůli stručnosti budeme nadále užívat termínu těžiště.} $C$ a
osu $z$ splývající s~osou symetrie levitátoru, dále pak hmotností $m$ a
hlavními centrálními momenty setrvačnosti $J_{\parallel}$ vzhledem k~rotační
ose levitátoru (osa $z$) a $J_{\perp}$ vzhledem k~libovolné ose kolmé
k~rotační ose levitátoru (např. osy $x, y$, které kvůli přehlednosti nejsou
znázorněny) a procházející těžištěm. Levitátor je umístěn v~homogenním
gravitačním poli s~komponentami $\vec{g}=(0,0,-g)$ v~laboratorním inerciálním
systému souřadnic $(X,Y,Z)$, kde $g>0$, a v~magnetickém poli generovaném výše
popsaným toroidálním bázovým magnetem. Osa $z$ svírá s~osou $Z$ nutační úhel
$\theta$. Písmeny S a J jsou označeny \uv{severní} a \uv{jižní} pól obou
magnetů; levitátor je zachycen v~poloze, kdy jsou k~sobě přivráceny
stejnojmenné magnetické póly a právě v~této situaci bude tedy docházet
k~žádoucímu odpuzování a nadnášení.

Levitátor~-- jako tuhé těleso bez restrikce na polohu jakéhokoli svého bodu~--
má šest stupňů volnosti
$(R\equiv\sqrt{X^{2}+Y^{2}},\Phi,Z,\varphi,\theta,\psi)$, z~nichž orbitální
úhel těžiště $\Phi$, precesní úhel $\varphi$ a úhel vlastní rotace $\psi$
nejsou na obrázku~\ref{UkazSchemaLevitronu} kvůli přehlednosti explicitně
vyznačeny; v~tomto směru odkazujeme na standardní učebnice analytické
mechaniky, např. na známou knihu Goldsteina a
spoluautorů~\cite{Gol-Poo-Saf:2002:ClassMech:3}.

Při výpočtech zanedbáváme skutečnost, že magnet umístěný v~levitátoru má
určité konečné rozměry a nahrazujeme jej bodovým dipólem, dále pak zanedbáváme
šířku prstence bázového permanentního magnetu budicího magnetické pole, a
nahrazujeme jej lineárním kružnicovým magnetem infinitezimální šířky a
výšky. Tato zanedbání, jak uvidíme, nic nemění na kvalitativním způsobu pohybu
levitátoru. Bylo by ale podnětné zkusit do budoucna složitější výpočty beroucí
v~úvahu rozměry báze i~levitátoru. Pro zpřehlednění vzorců elektromagnetismu
používáme místo soustavy jednotek SI soustavu
CGS~\cite{Pur-Mor:2013:ElectMagn:3}.

Za pozornost stojí, že magnetický dipól (v~praxi malý permanentní magnet) je
umístěn mimo těžiště $C$ (avšak stále na ose symetrie levitátoru); jejich
vzdálenost je na obrázku~\ref{UkazSchemaLevitronu} označena~$s$, přičemž ve
znázorněné \uv{odpuzovací} poloze je těžiště pod dipólem. Tato představa je
užitečná, protože teoreticky by se dalo dostatečnou volbou vzdálenosti~$s$
dosáhnout kompenzace a magnetického silového momentu $\vec{\mu}\times\vec{B}$
snažícího se překlopit levitátor gravitačním silovým momentem, kdy by
levitátor fungoval v~podstatě jako kyvadlo s~centrem magnetické interakce
v~dipólu a centrem gravitační interakce v~těžišti. Tento způsob stabilizace by
však, jak zakrátko uvidíme, nefungoval, a pro rovnice numerického modelu se
nakonec stejně klade $s=0$, což odpovídá reálnému zařízení.

Z~důvodu rotační symetrie Levitronu bude rovnovážná poloha levitátoru ležet na
ose bázového magnetu $Z$, a osa levitátoru $z$ s~ní bude antiparalelní. Pro
určení rovnovážné polohy proto můžeme potlačit pět stupňů volnosti buď fixací
jejich hodnot $(R=0,\theta=\pi)$, kdy jsou hodnoty $\Phi$, $\varphi$, $\psi$,
jakož i~případná nenulová hodnota $s$, irelevantní. Potenciální
energie~(\ref{PotEn}) se zjednoduší na tvar\footnote{Pro existenci rovnováhy
  musí mířit magnetická indukce v~rovnovážném bodě vzhůru, odkud plyne změna
  znaménka u~magnetického členu.}
\begin{equation}
  U = mgZ - \vec{\mu}\cdot\vec{B} = mgZ + \mu B_{Z}\,.              \label{r3_1}
\end{equation}
V~rovnováze musí síla působící ve směru $Z$ na levitátor vymizet,
\begin{equation}
  F_{Z}=-\pder{U}{Z}=-mg-\mu\pder{B_{Z}}{Z}=0\,.                     \label{r3_2}
\end{equation}
Podrobnější diskuse pomocí numerického
modelování~\cite{Hle:2011:SpinStabMagLev:WTC2011} vede k~závěru, existují
čtyři rovnovážné polohy, jak je vidět na obrázku~\ref{EqPoints}:
\begin{itemize}
\item Dvě \uv{vznášející se}, které mají $Z>0$, $\vec{\mu}$ míří dolů, báze a
  levitátor se odpuzují. Tyto rovnovážné polohy jsou dány
  rovnicí~(\ref{r3_2}).
\item Dvě \uv{zavěšené}, které mají $Z<0$, $\vec{\mu}$ míří nahoru, báze a
  levitátor se přitahují. Tyto rovnovážné polohy jsou dány modifikací rovnice
  rovnicí~(\ref{r3_2}) pro paralelní indukci báze a $\vec{\mu}$ projevující se
  změnou znaménka u~magnetického členu,
\begin{equation}
  F_{Z}=-\pder{U}{Z}=-mg+\mu\pder{B_{Z}}{Z}=0\,.                 \label{r3_2susp}
\end{equation}
 Těmito dvěma body se dále nebudeme zabývat, jejich vlastnosti kopírují
 vlastnosti \uv{vznášející se} dvojice rovnovážných bodů.
\end{itemize}
Magnetická indukce báze je úměrná jejímu celkovému magnetickému
dipólovému momentu $M$. Z~rovnic~(\ref{r3_2}) a~(\ref{r3_2susp}) je
pak zřejmé, že aby vůbec existovalo řešení a rovnovážné body, musí být
\uv{vazební koeficient} definovaný součinem
\begin{equation}
  \alpha \equiv \mu M                                          \label{BoundCoef}
\end{equation}
velikostí dipólových momentů levitátoru $\mu$ a báze $M$ větší než určité
limitní hodnota. Není-li tato podmínka splněna, levitace nenastane.

\begin{figure}[t]
\begin{center}

\includegraphics[width=\linewidth]{EquilibriumPoints}
\end{center}
\caption[Rovnovážné body Levitronu]{\label{EqPoints}Rovnovážné body Levitronu.
  Osa $Z$ je zde položena vodorovně, směr vzhůru je orientován
  doprava. \uv{Vznášející se} rovnovážné body jsou označeny~3, 4,
  \uv{zavěšené} rovnovážné body jsou označeny~1, 2. Tyto dvojice rovnovážných
  bodů se nacházejí vždy nad a pod inflexním bodem magnetické indukce na ose
  $Z$ (označeno \uv{infl}), a vždy v~místech, kde magnetická indukce
  slábne. V~bodech \uv{max} a \uv{$B=0$} je magnetická indukce maximální
  resp. nulová.  Převzato z~\cite{Hle:2011:SpinStabMagLev:WTC2011}.}
\end{figure}

Jednoduchým výpočtem (pro podrobnosti odkazujeme opět
na~\cite{Hle:2011:SpinStabMagLev:WTC2011}) můžeme ověřit, že pro nejhořejší
rovnovážný 4 bod je
\begin{equation}
  \pdder{U}{Z} > 0\,,                                             \label{dUZgt0}
\end{equation}
a tedy levitátor je vertikálně stabilní, zatímco pro rovnovážný bod 3 platí
opačná nerovnost
\begin{equation}
  \pdder{U}{Z} < 0\,,                                             \label{dUZlt0}
\end{equation}
takže levitátor je vertikálně labilní. Pro realizaci levitace tak připadá do
úvahy jedině nejhořejší rovnovážný bod~4.

Vraťme se ještě k~předchozí úvaze o~možnosti stabilizace levitátoru pomocí
posunu těžiště vůči magnetickému dipólu. Moment tíhové síly by zajišťoval
antiparalelní orientaci dipólu levitátoru, a~při presenci vhodného tlumicího
média (vzduch) by zajišťoval jeho exaktně nebo téměř vertikální
pozici. Rovnice~(\ref{r3_1}) by tedy (aspoň s~přesností do členů 2.~řádu)
platila i~při malých laterálních výchylkách, a díky harmonicitě polí $V$ a
$B_{Z}$, a~tudíž i~pole $U$, by v~rovnovážném bodě~4 v~důsledku~(\ref{dUZgt0})
a rotační symetrie muselo platit
\begin{equation}
  \pdder{U}{X} = \pdder{U}{Y} < 0\,,                             \label{dUXYlt0}
\end{equation}
takže poloha levitátoru by byla laterálně (horizontálně)
labilní.\footnote{Analogicky by se dalo dokázat, že v~rovnovážném bodě~3 by
  byl levitátor laterálně sice stabilní, ale vertikálně labilní, jak
  ukazuje~(\ref{dUZlt0}).} Levitátor by tedy měl tendenci v~tomto směru
\uv{sklouznout} a \uv{uletět}. Je zřejmé, že statický způsob zabránění
překlopení levitátoru nefunguje, neboť ponechává labilitu v~horizontálním
směru, plně v~souladu s~Earnshawovým teorémem.

Nyní výše presentované úvahy zpřesníme, čímž se zároveň připravíme na
sestavení matematického modelu levitátoru. Budeme zhruba sledovat postup
podle~\cite{Hab:1999:MgrDP:ProbUlFyzPrakt}. Položme počátek soustavy $(X,Y,Z)$
do místa, kde magnetická síla kompenzuje tíhu (takže levitátor levituje
v~okolí bodu $R=0$, $Z=0$) a rozviňme magnetickou indukci charakterizovanou
radiální a axiální komponentou $(B_{R}, B_{Z})$ v~okolí tohoto bodu do
mocninné řady s~přesností do kvadratických členů,
\begin{align}
  B_{R}(R,Z)=&B_{R}(0,0)+\pder{B_{R}}{R}(0,0)R+\pder{B_{R}}{Z}(0,0)Z\notag\\
  &+\frac{1}{2}\pdder{B_{R}}{R}(0,0)R^{2}+\frac{\partial^{2} B_{R}}
  {\partial{R}\partial{Z}}(0,0)RZ
  +\frac{1}{2}\pdder{B_{R}}{Z}(0,0)Z^{2}\,,                         \label{r3_11}
\end{align}
\begin{align}
  B_{Z}(R,Z)=&B_{Z}(0,0)+\pder{B_{Z}}{R}(0,0)R+\pder{B_{Z}}{Z}(0,0)Z\notag\\
  &+\frac{1}{2}\pdder{B_{Z}}{R}(0,0)R^{2}+\frac{\partial^{2} B_{Z}}
  {\partial{R}\partial{Z}}(0,0)RZ+\frac{1}{2}
  \pdder{B_{Z}}{Z}(0,0)Z^{2}\,.                                     \label{r3_12}
\end{align}
Zřejmě díky symetrii platí
\begin{equation}
  B_{R}(0,0)=\pder{B_{R}}{Z}(0,0)= \pdder{B_{R}}{Z}(0,0)=0         \label{axsym1}
\end{equation}
a také
\begin{equation}
  \pder{B_{Z}}{R}(0,0)=\pdderx{B_{Z}}{R}{Z}(0,0)=0\,.             \label{axsym2}
\end{equation}
Dále, vektorové pole $\vec{B}$ splňuje rovnice nezřídlovosti a nevírovosti
\begin{equation}
  \nabla\cdot\vec{B}=\pder{B_{Z}}{Z}+\frac{1}{R}\pder{}{R}(R B_{R})=0\,,\quad
  \left(\nabla\times\vec{B}\right)_{\Phi}=\pder{B_{R}}{Z}-\pder{B_{Z}}{R}=
  0\,,                                                             \label{r3_13}
\end{equation}
ostatní komponenty rotace $\nabla\times\vec{B}$ jsou identicky nulové díky
$B_{\Phi}\equiv 0$ a axiální symetrii magnetického pole báze $\vec{B}$.

Z~rovnic~(\ref{r3_11})--(\ref{r3_13}) plynou následující rozvoje s~přesností
do kvadratických členů~\cite{Hab:1999:MgrDP:ProbUlFyzPrakt}
\begin{align}
  B_{R}(R,Z)&=-\tfrac{1}{2}S_{0}R-K_{0}RZ\,,                       \label{r3_14}\\
  B_{Z}(R,Z)&=B_{0}+S_{0}Z+K_{0}Z^{2}-\tfrac{1}{2}K_{0}R^{2}\,,        \label{r3_15}
\end{align}
v~nichž je použito zkrácených označení
\begin{align}
  B_{0}&\equiv B_{Z}(0,0)\,, \notag\\
  S_{0}&\equiv \pder{B_{Z}}{Z}(0,0)=-2\pder{B_{R}}{R}(0,0)\,,      \label{r3_16}\\
  K_{0}&\equiv \frac{1}{2}\frac{\partial^{2}B_{Z}}{\partial{Z^{2}}}(0,0)=
  -\frac{\partial^{2}B_{R}}{\partial{R}\partial{Z}}(0,0)=
  -\frac{\partial^{2}B_{Z}}{\partial{R^{2}}}(0,0)\,.
\end{align}
Kdyby byla orientace magnetického momentu levitátoru stále svislá, jak jsme
již diskutovali výše, potenciální energie by se dala v~okolí rovnovážného bodu
vyjádřit s~přesností do kvadratických členů
\begin{equation}
  U= mgZ-\vec{\mu}\cdot\vec{B} =\mu B_{0}+(mg + \mu S_{0})Z
    +\mu K_{0}Z^{2}-\tfrac{1}{2}\mu K_{0}R^{2}\,.                    \label{r3_17}
\end{equation}
Pro vykompenzování tíhy levitátoru se musí anulovat koeficient u~$Z$. Aby však
rovnováha byla stabilní, musela by kvadratická část potenciální
energie~(\ref{r3_17}) být pozitivně definitní, což, jak je vidět, díky
poslednímu členu neplatí. Tento nestabilní režim je ilustrován na
obrázku~\ref{UnstabMode}.

\begin{figure}[t]
\begin{minipage}[b]{0.63\linewidth}
\begin{center}

\includegraphics[width=\linewidth]{InteractModelUnstab3}
\end{center}
\end{minipage}\hfill%
\begin{minipage}[b]{0.33\linewidth}
\begin{center}

\includegraphics[width=\linewidth]{InteractModelUnstab2}
\end{center}
\end{minipage}
\caption[Ilustrace nestabilního režimu Levitronu]{\label{UnstabMode}Ilustrace
  principu nestability v~nejvyšším rovnovážném bodě~4, kdy je osa levitátoru
  při výchylkách fixována ve vertikálním směru. Průběh potenciální energie
  v~libovolné rovině obsahující osu $Z$ v~okolí rovnovážného bodu~4 je
  znázorněn na barevném panelu uprostřed pomocí barevné mapy (vzrůst je
  reprezentován barvami duhy od červené po fialovou). Rovnovážná poloha je
  reprezentována bílým bodem, v~jehož okolí je patrné zřetelně sedlový průběh
  klesající horizontálně (laterálně) a rostoucí vertikálně. Převzato
  z~\cite{Hle:2011:SpinStabMagLev:WTC2011}.}
\end{figure}

Je tedy zřejmé, že rotace levitátoru má kromě stabilizace vůči překlopení
magnetického momentu levitátoru $\vec{\mu}$ do směru $\vec{B}$ ještě jiný,
velmi delikátní úkol~-- musí zajistit pozitivní definitnost kvadratické části
potenciální energie~(\ref{r3_17}) a tím i~prostorovou stabilitu
levitátoru. Jak toho lze dosáhnout? Na základě úvah sekce~\ref{BeyondE} na
straně~\pageref{BeyondE} zřejmě levitátor a jeho magnetický dipól $\vec{\mu}$
nemůže zůstat vertikální, ale musí se natáčet do směru lokálního magnetického
pole. Přesněji řečeno, magnetický moment $\vec{\mu}$ bude konat precesní pohyb
okolo směru lokální magnetické siločáry.  Pak místo výrazu
$\vec{\mu}\cdot\vec{B}=-\mu B_{Z}$ nastoupí (alespoň s~přibližnou platností)
výraz $\vec{\mu}\cdot\vec{B}\cong -\mu B$, kde $B$ označuje velikost
magnetické indukce $\left(B_{R}^{2}+B_{Z}^{2}\right)^{1/2}$. V~takovém případě
dostaneme (s~přesností do kvadratických členů) výraz~(\ref{r3_17})
modifikovaný na tvar~\cite{Hab:1999:MgrDP:ProbUlFyzPrakt}
\begin{align}
  U = \mu B_{0}&+(\mu S_{0} +mg)Z\notag\\
    &+\mu K_{0} \left(\frac{S^{2}_{0}}{2B_{0}K_{0}}+1\right)Z^{2}
    +\frac{1}{2}\mu K_{0}\left(\frac{S^{2}_{0}}{4B_{0}K_{0}}
      -1\right)R^{2}\,.                                            \label{r3_18}
\end{align}
Magnetická indukce na ose $Z$ (nyní vrátíme počátek soustavy $(X,Y,Z)$ opět do
středu bázového permanentního magnetu) je
\begin{equation}
  B_{Z}(0,Z)=\int
  \limits_{\Phi'=0}^{2\pi}\left[\frac{3MZ^{2}}{(Z^{2}+a^{2})^{\frac{5}{2}}}-
  \frac{M}{(Z^{2}+a^{2})^{\frac{3}{2}}}\right]\frac{\dif{\Phi'}}{2\pi}=
  M\frac{2Z^{2}-a^{2}}{(a^{2}+Z^{2})^{\frac{5}{2}}}\,,                  \label{r3_19}
\end{equation}
a její parciální derivace podle $Z$
\begin{equation}
  S = \pder{B_{Z}}{Z}
    = 3MZ\frac{3a^{2}-2Z^{2}}{(a^{2}+Z^{2})^{\frac{7}{2}}}\,.           \label{r3_20}
\end{equation}
Podmínka anulování koeficientu $\mu S+mg$ a tedy rovnováhy je
\begin{equation}
  3\mu M Z~\frac{2Z^{2}-3a^{2}}{(a^{2}+Z^{2})^{\frac{7}{2}}}=mg\,,      \label{r3_21}
\end{equation}
z~níž pro zadané hodnoty $\mu$, $M$, $m$ a $g$ numericky vypočteme levitační
výšku $Z_{0}$.  Potenciální energie je tedy (nepodstatný konstantní člen
$\mu B_{0}$ vynecháme)
\begin{equation}
  U = \mu K_{0} \left(\frac{S^{2}_{0}}{2B_{0}K_{0}}+1\right)(Z-Z_{0})^{2}
  +\frac{1}{2}\mu K_{0}
  \left( \frac{S^{2}_{0}}{4B_{0}K_{0}}-1\right)R^{2}\,,               \label{r3_22}
\end{equation}
nebo ekvivalentně
\begin{equation}
  U = \mu K_{0}
  \left(\frac{m^{2}g^{2}}{2B_{0}K_{0}\mu^{2}}+1\right)(Z-Z_{0})^{2}+
  \frac{1}{2}\mu K_{0}
    \left(\frac{m^{2}g^{2}}{4B_{0}K_{0}\mu^{2}}-1\right)R^{2}\,.      \label{r3_23}
\end{equation}
přičemž $K_{0}=\partial^{2}B_{Z}/\partial Z^{2}$ a $B_{0}$ se bere v~bodě
$R=0$, $Z=Z_{0}$, kde $Z_{0}$ je řešením
rovnice~(\ref{r3_21}). Rovnice~(\ref{r3_21}) může mít v~oblasti $Z>0$ dvě
řešení~-- je nutno vybrat to z~nich, pro které jsou koeficienty
u~kvadratických členů v~(\ref{r3_22}) nebo~(\ref{r3_23}) kladné. Parametry
levitátoru tedy musí být voleny tak, aby
\begin{equation}
  K_{0}>0\,,\quad \frac{m^{2}g^{2}}{4B_{0}K_{0}\mu^{2}}>1\,.           \label{r3_24}
\end{equation}
\begin{figure}[t]
\begin{minipage}[b]{0.63\linewidth}
\begin{center}

\includegraphics[width=\linewidth]{InteractModelStab3}
\end{center}
\end{minipage}\hfill%
\begin{minipage}[b]{0.33\linewidth}
\begin{center}

\includegraphics[width=\linewidth]{InteractModelStab2}
\end{center}
\end{minipage}
\caption[Ilustrace stabilního režimu Levitronu]{\label{StabMode}Ilustrace
  princip stabilního režimu v~nejvyšším rovnovážném bodě~4, kdy osa levitátoru
  při výchylkách sleduje směr siločar pole $\vec{B}$. Průběh potenciální
  energie v~libovolné rovině obsahující osu $Z$ v~okolí rovnovážného bodu~4 je
  znázorněn na barevném panelu uprostřed pomocí barevné mapy (vzrůst je
  reprezentován barvami duhy od červené po fialovou). Rovnovážná poloha je
  reprezentována bílým bodem, v~jehož okolí je patrné prostorové lokální
  minimum zajišťující stabilitu. Převzato
  z~\cite{Hle:2011:SpinStabMagLev:WTC2011}.}
\end{figure}
V~tom případě je levitátor polapen v~prostorové potenciálové jámě a jeho
těžiště osciluje v~horizontální rovině s~periodou
\begin{equation}
  T_{\mathrm{R}}\sim 2\pi\sqrt{\frac{m}{2\gamma\,_{\mathrm{R}}}}\,,    \label{r3_25}
\end{equation}
kde $\gamma\,_{\mathrm{R}}$ je koeficient u~$R^2$ v~rovnici~(\ref{r3_22})
nebo~(\ref{r3_23}), analogicky ve vertikálním směru těžiště osciluje
s~periodou
\begin{equation}
  T_{\mathrm{Z}}\sim 2\pi\sqrt{\frac{m}{2\gamma\,_{\mathrm{Z}}}}\,,    \label{r3_26}
\end{equation}
kde $\gamma\,_{\mathrm{Z}}$ je koeficient u~$Z^2$ v~(\ref{r3_22})
nebo~(\ref{r3_23}). Dále, předpokládáme-li, že rotace levitátoru je dostatečně
rychlá, koná levitátor precesní pohyb kolem lokální siločáry s~periodou
\begin{equation}
  T_{\mathrm{prec}} \sim 2\pi\frac{J_\parallel\,\Omega_0}{\mu B_0} \ll
  T_{\mathrm{rot}}= \frac{2\pi}{\Omega_{0}}                           \label{r3_27}
\end{equation}
(viz~\cite{Gol-Poo-Saf:2002:ClassMech:3}), kde $\Omega_0\equiv\loder{\psi}{t}$
je úhlová rychlost rotace levitátoru. Jestliže platí adiabatická podmínka
\begin{equation}
  T_{\mathrm{prec}} \ll T_{\mathrm{R}}\,,                                \label{r3_28}
\end{equation}
stačí při výchylkách, při nichž levitátor prochází oblastmi s~různým sklonem
magnetických siločar, precesní osa sledovat změnu jejich směru, a tak vytvářet
v~horizontálním směru potenciálovou jámu. Odtud plyne zajímavý závěr: rotace
levitátoru nesmí být příliš rychlá, neboť pak je precesní perioda dlouhá a
stane se srovnatelnou nebo dokonce větší než $T_{\mathrm{R}}$, což vede ke
ztrátě stability. Rotace levitátoru má samozřejmě i~dolní mez: při poklesu
otáček pod ni přestane rotace účinně bránit překlopení $\vec{\mu}$ do směru
souhlasně s~$\vec{B}$. Fungování stabilního režimu je ilustrováno na
obrázku~\ref{StabMode}.

\endinput

%%
%% End of file `x_levit.tex'.
